\section*{Reproducibility and AI-use statement}
The public repository is \url{https://github.com/SzeChunYiu/Orion}. It maintains frozen benchmarks, content-addressed execution receipts, exact-subject CI, negative-history records, release manifests and editable paper/figure sources. The V2.1 engineering figures are rendered with deterministic Matplotlib code from \path{research/MINI_RESEARCH_DEMO_043_RECEIPT.json}, \path{research/MINI_RESEARCH_METROLOGY_044_RECEIPT.json} and \path{research/MINI_ARCHIVE_STORAGE_044_RECEIPT.json}; CI checks the plotted source coordinates and exports PDF, editable SVG, PNG preview and machine-readable figure source data. Structural arrows are intentionally unlabelled in the paper schematics, and quantitative arrow callouts are forbidden by repository tests. The first matched-study attempt is preserved separately as an executable freeze packet, a fail-closed preflight receipt with zero evaluated task records and a non-scientific provider-connectivity receipt. These are readiness artifacts, not model-performance data. Language models are research, coding and drafting tools rather than authors; their proposals do not possess canonical scientific authority. Same-context role simulations and the present internal Nature-style review rounds are not counted as independent peer review. The later journal-results manuscript will bind exact data/transformation identities, model-run receipts and all quantitative figure source data.

\appendix
\subsection{Reproducibility as a chain of identities}

Reproducibility is not one binary property. Computational reproducibility asks whether the same code, environment and data can regenerate the result. Analytical reproducibility asks whether the transformation from data to claim is specified enough to inspect. Experimental reproducibility asks whether a new realization of the scientific process yields compatible evidence. Conceptual replication asks whether a result survives a different operationalization of the same underlying question. These levels can disagree.

The repository infrastructure mainly addresses the first two. Exact subject SHAs, content hashes, deterministic receipts, dependency declarations and publication manifests allow a reader to identify the artifact that generated a figure or software claim. Established reproducible-research guidance motivates this bookkeeping \cite{sandve2013,wilson2017}, and broader reproducibility reform emphasizes transparency across design, analysis and reporting \cite{munafo2017}. The framework does not claim that such bookkeeping produces independent replication.

This hierarchy explains why artifact identity is printed inside the manuscript. If a software contract changes after a paper is rendered, the old PDF still describes the old subject. If the release process instead prints ``latest'' or an unpinned package version, future readers cannot know which behavior was tested. Content addressing converts that ambiguity into an immutable reference.

Data and model availability can still limit reproducibility. Proprietary APIs can change, source websites can disappear, and external models can be revised behind stable product names. When exact replay cannot be guaranteed, the receipt records the strongest available identity, raw outputs where permitted, timestamps, settings and hashes of locally controlled inputs. The residual limitation is reported explicitly rather than masked by a reproducibility label.

AI-in-the-loop scholarship is also moving toward executable provenance. F(AI)2R, for example, treats the paper-production process itself as an auditable provenance graph and uses continuous integration to check graph conformance \cite{krebs2026fair}. The important lesson for Orion is not that CI can certify scientific truth. It is that manuscript text, code, evidence and verification acts can be made addressable enough that later reviewers can ask which artifact authorized a sentence. The present release follows the same direction by binding generated engineering claims to exact source subjects and by separating same-context manuscript review from independent scientific replication.

Independent replication then sits above the artifact layer. A fresh group using the published method and independently acquired evidence can test whether the scientific effect transports. That event would upgrade a distinct authority coordinate. The current paper has not earned that coordinate and says so.

\subsection{Publication artifact identity and chaptered source}

The public paper package is generated as a chaptered \LaTeX{} project. Each major top-level section is stored in its own source file and assembled by a small \texttt{main.tex}; quantitative figure sources remain receipt-bound; and a source manifest records hashes for the assembled paper and section files. The compiled PDF is treated as a derivative release artifact rather than as the primary editable record.

The build identity deliberately distinguishes the \emph{scientific software subject} from the later commit that stores rendered artifacts. The manuscript prints the exact implementation SHA evaluated by the build. CI runs the software tests, regenerates receipt-bound figures, stages the paper from that exact subject, compiles with blocking LaTeX warnings enabled, inspects PDF metadata and embedded fonts, renders every page, and only then may store the PDF alongside the TeX package. This avoids a recursive identity problem in which committing the PDF changes the source SHA that the PDF itself claims to describe.

Visual inspection remains necessary even when LaTeX exits successfully. Broken glyphs, clipped panels, unreadable figure labels and page collisions are artifact failures that a source-level unit test may not detect. The release process therefore renders all pages to images and treats the rendered pages, not the source alone, as the authoritative layout check.
